\section{Methods}

\subsection{Targets, ligands and splits}

We used the DOCKSTRING dataset, which holds AutoDock Vina 1.1.2 scores of 260,155 drug-like ligands docked into 58 prepared receptors.\cite{dockstring,vina} Targets were grouped into 22 protein families, and entire families were assigned to one role: 39 training targets, 10 development targets and 9 test targets. Ligands were grouped by Bemis-Murcko scaffold and split into a training set of 156,095 molecules and development, two calibration and test sets of 26,015 molecules each, with no scaffold group in two sets. The existing networks are trained on the docking scores of the \kpNTargets{} training targets with 40{,}000 training ligands per target, and the ten development targets serve only for the choice of their checkpoint. The cross-fitted networks are trained on the docking scores of \kfFoldTrainMin{} to \kfFoldTrainMax{} of the \kfNTargets{} training and development targets, depending on the fold, with 40{,}000 training ligands per target, so that the development targets are part of their training data. Both groups screen new ligands against the known pockets of their training data. A hit is a ligand in the lowest 1\% of the scores of a pocket within the ligand set that is evaluated, and the evaluation ligands are 10,000 development ligands that were not used in training.

\subsection{Networks}

The pocket is a residue graph in which each node carries amino-acid identity, mean partial charge, the distribution of AutoDock atom types and distance features, and edges connect the 16 nearest residues with radial basis features of the distance.\cite{gilmer2017} A three-layer message-passing encoder turns this graph into residue embeddings. The ligand encoder is a three-layer message-passing network over the two-dimensional atom graph; neither encoder sees coordinates of the ligand or a docking pose. The three networks share these encoders (dimension 128) and differ in how they combine the ligand with the pocket. The dual encoder (\paramsDual{} parameters) scores a product of a ligand vector and a pooled pocket vector, so that the score of a screen against many pockets is a matrix product. The pooled-concatenation network (\paramsPc{} parameters) concatenates the pooled ligand vector and the pooled pocket vector and reads them with a multilayer perceptron. Both keep the ligand embedding independent of the pocket and are Jev-like. The cross-attention network (\paramsXa{} parameters) passes the ligand atom states through two cross-attention layers with four heads in which the ligand atoms attend to the residues of the pocket, and reads the pooled result with a multilayer perceptron, so that the representation of a ligand depends on the pocket and every pair needs the interaction layers. The constant-pocket network is the dual encoder trained with one synthetic pocket graph for every target. Every network has a hit-probability readout, which ranks the ligands, and a score readout.

\subsection{Training}

Every network was trained with AdamW (learning rate $3\times10^{-4}$, weight decay $10^{-4}$) on batches of 256 ligands of one training pocket, with the binary cross-entropy of the top-1\% hits of the training scores of the pocket plus 0.2 times a Huber loss on standardized scores, for at most 60{,}000 updates or 10 epochs, and the hit-probability readout was evaluated every 6{,}094 updates, with training stopped after three evaluations without improvement. Two groups of networks were trained. The cross-fitted networks were trained in five folds of whole protein families (the kinases; the nuclear receptors; the class A GPCRs, CYP450 and MAO-B; the aspartyl proteases, metalloproteases, serine proteases S1, NOS1 and PARP1; and PDE5A, PTGS2 and PTPN1) over the 49 training and development targets: a network of a fold is trained on the targets outside the fold with 40{,}000 training ligands per target, and its checkpoint is the evaluation point with the highest macro recall at the 10\% budget on six of these targets against 10{,}000 calibration ligands. This is the choice of a checkpoint on seen pockets. The dual encoder, the pooled network and the cross-attention network were trained with the real pockets and with one constant synthetic pocket for every target, with seed 11, and the dual encoder and the cross-attention network were also trained with a longer recipe (learning rate $10^{-4}$, up to 150{,}000 updates or 30 epochs, patience of eight evaluations). The existing networks come from a preceding analysis of the author (Supporting Information): the dual encoders and the constant-pocket networks, and the pooled and cross-attention networks that were trained with the same recipe on the 39 training targets, with three seeds each, and registered as addendum 8 before the first run; their checkpoint is the evaluation point with the highest macro recall on the ten development targets against 10{,}000 development ligands, the choice of a checkpoint on new targets. The recipe was built for the dual encoders and was not tuned for the other networks.

\subsection{Evaluation on known pockets}

Every network ranked the 10{,}000 evaluation ligands for each pocket of its training data with its hit-probability readout. The evaluation ligands are 10{,}000 development ligands that were not used in training and not in the choice of the checkpoint of the cross-fitted networks; the existing networks chose their checkpoint with the development targets against the same ligands. The top-1\% recall at a budget is the fraction of the hits of a pocket that lie in the best 1, 5, 10 or 20\% of its ranking. The recall of a pocket is the mean over the seeds (existing networks) or over the folds that train on the pocket (cross-fitted networks, four of five), and values are macro means over the pockets. The 95\% intervals are percentile intervals from 5{,}000 bootstrap resamples of the pockets, and the paired difference of two networks is the difference of their per-pocket recalls resampled in the same way. The two evaluations were registered as addenda 10 and 11 before they were computed.

\subsection{Cost at the full scale}

The cost was measured at the full scale on one \tmAGpu{} rented from RunPod (PyTorch 2.4.1, CUDA 12.4), for all \tmALig{} ligands of the DOCKSTRING collection screened from SMILES against 1, 5, 20 and 57 wild-type pockets, and against 1{,}000 pockets obtained by using the 57 pockets cyclically (every copy is encoded and scored again). The stages are the featurization of the SMILES on \tmAProcs{} CPU processes, the collation of the batches, the embedding of the ligands on the GPU, the pocket states of all pockets, the scoring of all pairs and the selection of the best 10\% of the ligands of every pocket with \texttt{torch.topk}; the model time is the sum of the GPU stages. Each GPU stage was synchronized and timed twice, and the shorter time is reported. The dual encoder scores all pairs with one matrix product on the cached ligand and pocket vectors. The pooled network decomposes the first layer of its readout into a ligand term and a pocket term that are computed once and combined for each pair. The cross-attention network keeps the pocket-independent ligand atom states of the whole library on the GPU without padding, passes them through the interaction layers for each pocket with fused attention, and was timed in float32 and in float16 for the interaction layers and the readout, and its original padded implementation was timed on a subset for reference. Every implementation was checked against the module forward on 2{,}000 ligands and five pockets (rank correlation of at least 0.9999 in float32 and 0.999 in float16). The Supporting Information gives every stage.

Uni-Dock 1.2.0 was timed on one RTX 4090 rented from RunPod (driver 570.169), the GPU class of the docking runs, with the Vina scoring function, one output pose and the DOCKSTRING box of each receptor. For three receptors (ABL1, MAOB and PTGS2) the wild-type receptor was docked against the first $n$ of 1,023 prepared ligands (\dockPrepOkFourK{} for the batches of 4,000), a random sample of the DOCKSTRING collection, in every combination of search mode (detail, balance and fast) and batch size (160, 256, 1,000 and 4,000 ligands per call, 160 for the detail mode only), and the rate is the number of ligands with a score divided by the wall time of the \texttt{unidock} process. The time that docking needs for a screen is the number of pairs divided by the median rate of the configuration with the highest rate. It leaves out the preparation of the ligand files (protonation, embedding and PDBQT conversion), which the networks do not need, whereas the time of the networks includes the featurization of the ligands from SMILES.

\subsection{Registered plans and disclosure}

The comparison with joint networks and the full-scale timing were registered as addenda 8 and 9 of the analysis plan before the respective runs. The restriction to known pockets (addendum 10) was written after most of the experiments of addendum 9 had finished, and the evaluation of the cross-fitted networks (addendum 11) was written after the outcomes of the evaluation of the existing networks were known, because the checkpoint rule of the existing networks selects for new targets. Both evaluations are reported. The plan with its hash ledger is part of the data archive. The dual encoders and the constant-pocket networks come from a preceding, unpublished analysis of the author, the pooled and cross-attention networks were trained in this study, and the Supporting Information lists every addendum, the reads of the sealed test labels and the analyses that followed the registered outcomes.
